\documentclass{article}
\usepackage[utf8]{inputenc}
\usepackage[T1]{fontenc}
\usepackage{geometry}
\usepackage{graphicx}
\usepackage{float}
\usepackage{hyperref}
\usepackage[english,indonesian]{babel}

\title{Optimasi Pemrosesan Template LaTeX Menggunakan Pendekatan Berbasis AI}
\author{Academic Support Agent}
\date{\today}

\begin{document}
\maketitle

\begin{abstract}

\end{abstract}

\chapter{PENDAHULUAN}


\chapter{TINJAUAN PUSTAKA}


% Original content (commented out):
% \documentclass[a4paper,11pt]{article}

% Import des fichiers de configuration
\input{0_config/01_packages}
\input{0_config/02_commands} 
\input{0_config/03_variables}

\setlength{\headheight}{24pt}
\addtolength{\topmargin}{-8pt}

% Construction de la couverture
\input{1_cover/11_covermain}

\begin{document}

% Contenu principal
\input{2_content/20_introduction}

\newpage
\input{2_content/21_part}

\newpage
\input{2_content/22_part}

\newpage
\input{2_content/23_part}

\newpage
\input{2_content/24_part}

\newpage
\input{2_content/28_conclusion}

% Bibliographie
\newpage 
\input{2_content/29_bibliography}

\end{document}

\end{document}